% Ad Quintum Fratrem 2.11 --- headnote
% ad-quintum-fratrem-02-11, 14 February 54 BC, written at Rome

\textit{Marcus to Quintus, written at Rome on 14 February 54 BC
(\textit{a. d. xvi K. Martias}). Quintus is now in Gaul on
Caesar's staff --- a posting his elder brother had taken pains
to broker --- and the daily reports from Rome serve as much to
keep his spirits up as to inform him. Quintus's last letter
had evidently included some piece of paradox he labelled
``black snow,'' and Marcus picks up the joke approvingly: the
brother is in good humour and ready to play, which delights
him. The line on Pompey and Caesar is the political ground
of the season: with Crassus heading east, Pompey hovering
over Italy, and Caesar absorbing Gaul, the triumvirate's grip
is by now total, and Cicero --- chastened by the events of
56 and 55 BC --- is openly Caesarean. ``He is in my heart,''
he tells Quintus, ``and I do not loosen the bond.''}

\textit{The middle sections are a hurried bulletin of city
news. Cicero's young protégé M. Caelius Rufus is to be
prosecuted again, this time it seems by the loathsome Pola
Servius and at the urging of the Clodian faction --- the same
network whose attack he had repelled, with Cicero's help, in
the \textit{Pro Caelio} two years before. The Tyrian embassy
against the Syrian tax-farmers brought a packed Senate, where
Gabinius (back from his Syrian governorship and shortly to
face his own trials) was savagely handled, and L. Aelius
Lamia, defending the equestrian jurors against Domitius
Ahenobarbus, scored a clean retort: ``\textit{We} give the
verdicts; \textit{you} pronounce the praises.'' The third
section turns to procedure --- the consul Appius Claudius
Pulcher reading the lex Pupia and lex Gabinia so as to keep
the Senate sitting daily for embassies through February,
which will likely shove the elections into March --- and the
fourth, the most affectionate, drops politics altogether for
the Greek historians Quintus has been reading on campaign:
Callisthenes, Philistus of Syracuse (``almost a Thucydides
in miniature''), and Dionysius. Marcus urges his brother to
take up history himself, promises a fresh dispatch on the
Lupercalia (15 February), and ends with a hug for the
nephew, the young Cicero.}
